% !PRINTED 13
\weberpage{13}

Setzt man dann
\weq{9}{z^2 = y,}
so ergibt sich
\weq{10}{\frac{2\,dz}{\sqrt{(z^2-\alpha)(z^2-\beta)}} = \frac{dy}{\sqrt{y(y-\alpha)(y-\beta)}}}
und man hat also durch die \gesperrt{quadratische Substitution} (9) ein elliptisches Differential erhalten, bei dem unter dem Wurzelzeichen eine Funktion dritten Grades mit reellen Wurzeln steht, das man nach \S\ 3 durch lineare Substitution auf die Normalform
\[
\frac{d\zeta}{\sqrt{\zeta(1-\zeta)(1-\varkappa^2\zeta)}}
\]
bringen kann, und darin können $\zeta$ und $\varkappa$ als positive echte Brüche angenommen werden.

Die \gesperrt{Legendresche} Normalform ergibt sich daraus, wenn man
\[
\zeta = \sin^2\varphi, \qquad d\zeta = 2\sin\varphi\cos\varphi\,d\varphi
\]
setzt:
\weq{11}{\frac{d\zeta}{\sqrt{\zeta(1-\zeta)(1-\varkappa^2\zeta)}} = \frac{2\,d\varphi}{\sqrt{1-\varkappa^2\sin^2\varphi}}.}

\begin{center}
\S\ 5. Die \gesperrt{Weierstrasssche} Normalform.
\end{center}

Eine andere Normalform des elliptischen Differentials hat \gesperrt{Weierstrass} seinen Untersuchungen zugrunde gelegt, nämlich die Form
\weq{1}{du = \frac{dz}{\sqrt{4z^3 - g_2 z - g_3}},}
in der $g_2$, $g_3$ Konstanten sind, die die erste und zweite Invariante des Differentials genannt werden. Das allgemeine elliptische Differential
\weq{2}{\frac{dx}{\sqrt{f(x)}},}
worin
\weq{3}{f(x) = a_0 x^4 + a_1 x^3 + a_2 x^2 + a_3 x + a_4}
ist, kann durch eine \gesperrt{lineare Transformation} auf die Weierstrasssche Normalform gebracht werden, wenn die Wurzeln $x_1$, $x_2$, $x_3$, $x_4$ von $f(x)$ bekannt sind. Am einfachsten geschieht dies auf folgende Weise.
